\subsection{Domain and Range of a Linear Function} 
To determine the implicit domain of a function, we ask:\makeblanknewf\lbr
For a linear function, we have two possible operations involving the variable:
\begin{itemize}
\item \makeblankfill{}
\item \makeblankfill{}
\end{itemize}
These are both defined for \makeblank, so the domain of a linear function is \makeblank, in notation, \makeblanksmall or \makeblank[1in].\lbr
To find the range of a function, we have two possible options.
\begin{enumerate}
\item Consider the possible \makeblank[1.5in] of the operations involved.
\item Consider the domain of the \makeblank.
\end{enumerate}
Since linear functions are \makeblank, we choose option \makeblanksmall. The inverse of a linear function is \makeblank[1in], so the domain of the of the inverse is \makeblanksmall.
\begin{fact}[Domain and Range of a Linear Function]
If $f$ is a linear function, then $\dom f=\R$ and $\rng f=\R$.
\end{fact}